\chapter{Literature Review}
\section{Literature Review}

\subsection{Plant Disease Detection Systems}
\begin{itemize}
    \item Previous research studies will show that machine learning and CNN-based models will be widely used for detecting plant diseases from leaf images \cite{1}.
    \item Many systems will use datasets such as PlantVillage, which will help in improving accuracy and early identification of crop diseases \cite{2}.
    \item These studies will indicate that automated disease detection will help farmers reduce crop losses and take preventive measures in time \cite{3}.
\end{itemize}

\subsection{Soil Analysis and Crop Recommendation Systems}
\begin{itemize}
    \item Literature on smart farming systems will suggest that soil parameter evaluation such as pH, nitrogen, phosphorus, and potassium levels will be essential for selecting suitable crops \cite{4}.
    \item Existing research will highlight that data-driven soil analysis will improve soil fertility management and increase agricultural productivity \cite{5}.
    \item Studies will also show that crop recommendation models will support farmers in making better cultivation decisions \cite{6}.
\end{itemize}
\newpage

\subsection{AI-Based Chatbots in Agriculture}
\begin{itemize}
    \item Prior works will demonstrate that chatbots and conversational systems will be used to provide agricultural guidance and query-based support to farmers \cite{7}.
    \item These systems will improve accessibility to information in rural and remote areas by offering instant responses and simplified communication \cite{8}.
    \item Literature will indicate that chatbots will enhance farmer awareness and reduce dependency on manual consultation \cite{9}.
\end{itemize}

\subsection{Digital Agricultural Marketplaces}
\begin{itemize}
    \item Existing platforms will show that e-commerce solutions in agriculture will enable farmers to purchase seeds, fertilizers, pesticides, and tools conveniently \cite{10}.
    \item Studies will reveal that online marketplaces will reduce middlemen involvement and provide better transparency in pricing and product availability \cite{11}.
    \item Research will also suggest that digital platforms will support cost-effective procurement and improve farmer accessibility to agricultural resources \cite{12}.
\end{itemize}

\newpage

\section{Existing System}

\subsection{Plant Disease Detection Systems}
\paragraph{}Existing plant disease detection systems use machine learning algorithms to identify diseases from plant leaf images. These systems provide accurate results and help farmers take preventive measures. However, most of these applications focus only on disease detection and lack integration with other agricultural services such as soil analysis or product availability.

\subsection{Soil Analysis and Crop Recommendation Systems}
\paragraph{}Soil analysis systems are designed to evaluate soil parameters such as pH level, nitrogen, phosphorus, and potassium content. Based on these values, they recommend suitable crops and fertilizers. Although these systems improve decision-making, they often require technical knowledge and are not user-friendly for farmers.

\subsection{Agricultural Chatbot Systems}
\paragraph{}Chatbot-based systems are developed to provide agricultural guidance and answer farmer queries. These systems improve accessibility to information and provide real-time assistance. However, they are generally limited to advisory services and do not integrate other functionalities such as disease detection or marketplace features.

\subsection{Online Agricultural Marketplaces}
\paragraph{}Digital marketplaces allow farmers to purchase seeds, fertilizers, and farming tools online. These platforms reduce dependency on middlemen and improve accessibility to products. However, they operate independently and do not provide analytical or advisory support to farmers.

\newpage

\section{Proposed System}

\subsection{Krushi}
\paragraph{}The proposed system “Krushi” will be a user-friendly and dynamic web application designed to assist farmers in multiple agricultural activities. The system will integrate various modules such as plant disease detection, soil analysis, AI-based chatbot, and an online agricultural shop into a single platform. Farmers will be able to upload plant images to detect diseases and receive preventive suggestions. They will also be able to input soil parameters and get recommendations for suitable crops and fertilizers.

\paragraph{}Additionally, the system will provide a chatbot interface through which users can ask agricultural queries and receive instant responses. The integrated marketplace will allow farmers to purchase seeds, fertilizers, and tools easily. The platform will be accessible from anywhere, enabling farmers to utilize its services efficiently. By combining multiple features into one system, Krushi will improve decision-making, reduce crop losses, and promote digital agriculture practices.